%=============================================================================
% SQMS PHASE I FULL PROPOSAL
% "Precision Multi-Frequency Quality Factor Measurements in SRF Cavities
%  as a Test of Scalar-Field Modifications to Electrodynamics"
%
% Prepared for submission to:
%   Superconducting Quantum Materials and Systems (SQMS) Center
%   Fermi National Accelerator Laboratory
%   Director: Anna Grassellino
%
% Funding request: $3.2M over 24 months
% DOE Office of Science, Quantum Information Science programme
%
% Wave 4, Iteration 10 | DFD Research Programme
% Building on: W3i7 SQMS Protocol, W3i8 Detection Ranking,
%              W3i9 Channel B Audit, MASTER_FINDINGS_CORPUS
% Date: 20 March 2026
%=============================================================================

\documentclass[11pt]{article}
\usepackage[utf8]{inputenc}
\usepackage{amsmath,amssymb,amsfonts,amsthm,mathtools}
\usepackage{geometry}
\geometry{margin=1in}
\usepackage{booktabs}
\usepackage{siunitx}
\usepackage{hyperref}
\usepackage{xcolor}
\usepackage{enumitem}
\usepackage{float}
\usepackage{array}
\usepackage{longtable}
\usepackage{graphicx}
\usepackage{fancyhdr}
\usepackage{titlesec}
\usepackage{tcolorbox}
\tcbuselibrary{skins,breakable}

% Custom commands
\newcommand{\dpsi}{\delta\psi}
\newcommand{\lam}{\lambda}
\newcommand{\lbone}{|\lambda - 1|}
\newcommand{\Qpsi}{Q_\psi}
\newcommand{\Meff}{M_{\rm eff}}
\newcommand{\ee}[1]{\times 10^{#1}}
\newcommand{\half}{\tfrac{1}{2}}
\newcommand{\kB}{k_{\mathrm{B}}}
\newcommand{\Heff}{H_n}
\newcommand{\muG}{\mu_0^{\rm gal}}
\newcommand{\GN}{G_N}
\newcommand{\SNR}{\mathrm{SNR}}
\newcommand{\Vcav}{V_{\rm cav}}

\definecolor{sqmsblue}{RGB}{0,51,102}
\definecolor{sqmsgold}{RGB}{204,153,0}

\hypersetup{
  colorlinks=true,
  linkcolor=sqmsblue,
  citecolor=sqmsblue,
  urlcolor=sqmsblue
}

\pagestyle{fancy}
\fancyhf{}
\lhead{\small SQMS Phase I Proposal --- Scalar-Field Q-Ratio Diagnostic}
\rhead{\small \thepage}
\renewcommand{\headrulewidth}{0.4pt}

\newtheorem{theorem}{Theorem}[section]
\newtheorem{proposition}[theorem]{Proposition}

\title{%
\vspace{-1cm}
{\large\textsc{Proposal to the SQMS Center, Fermilab}}\\[12pt]
{\LARGE\bfseries Precision Multi-Frequency Quality Factor\\
Measurements in SRF Cavities as a Test of\\
Scalar-Field Modifications to Electrodynamics}\\[12pt]
{\large\textit{The SQMS Phase~I Psi-Diagnostic Experiment}}\\[8pt]
{\normalsize Funding Request: \$3.2M over 24 months}
}

\author{%
\textbf{Principal Investigator}: [Name], Fermilab/SQMS\\
\textbf{Co-PI}: [Name], [Institution]\\
\textbf{Theory Lead}: G.T.\ Alcock\\[6pt]
\textit{In collaboration with the SQMS Center consortium}
}

\date{March 2026}

\begin{document}
\maketitle
\thispagestyle{empty}

\vfill
\begin{center}
\fbox{\parbox{0.9\textwidth}{
\textbf{Proposal Summary}\\[4pt]
We propose a 24-month precision measurement programme using four existing
SQMS nitrogen-doped niobium cavities to search for anomalous
frequency-dependent quality factor ratios that would constitute evidence
for a scalar field coupled to electromagnetism. The key diagnostic---the
ratio of residual losses at TM$_{011}$ (2.4~GHz) versus TM$_{010}$
(1.3~GHz)---yields 0.49 under the Density Field Dynamics (DFD)
scalar-field hypothesis and 3.41 under conventional BCS theory. A
measurement to 10\% precision decisively distinguishes these predictions.
The experiment leverages existing SQMS infrastructure, requires no new
cryogenic capital, and improves the current bound on the scalar coupling
parameter by a factor of 2--6$\times$. A null result is scientifically
valuable, establishing the world's most stringent constraint on
EM--scalar coupling from SRF cavity data. A positive result would be the
first laboratory evidence for a new fundamental field.\\[4pt]
\textbf{Total request}: \$3.2M \quad
\textbf{Duration}: 24 months \quad
\textbf{Key deliverable}: $\lbone < 7.8\ee{-10}$ (null) or detection (positive)
}}
\end{center}
\vfill

\newpage
\tableofcontents
\newpage

%=============================================================================
%=============================================================================
\part{Scientific Case}
%=============================================================================
%=============================================================================

%=============================================================================
\section{Executive Summary}
\label{sec:executive_summary}
%=============================================================================

Density Field Dynamics (DFD) is a scalar-field extension of gravity and
electrodynamics that replaces dark matter with a single scalar field
$\psi$ acting as a cosmological refractive index. The theory makes 14
testable predictions from a single free parameter, the electromagnetic
coupling constant $\lambda$. The current best laboratory constraint is
$\lbone < 1.56\ee{-9}$, derived from existing SQMS quality factor
measurements.

This proposal describes the \textbf{single most cost-effective
experiment} to either detect or further constrain the DFD scalar field.
We propose:

\begin{enumerate}[nosep]
\item A $2\times 2$ array of nitrogen-doped niobium TESLA cavities, two
  pairs in separate cryogenic environments, enabling systematic diagnosis
  through cross-correlation.
\item Multi-frequency quality factor measurements at TM$_{010}$
  (1.3~GHz), TM$_{011}$ ($\sim$2.4~GHz), and TE$_{011}$
  ($\sim$1.7~GHz), exploiting the dramatically different
  frequency scaling predicted by DFD versus BCS theory.
\item A 12-month science campaign with continuous $Q_0(t)$ monitoring,
  environmental regression, and blind analysis protocol.
\end{enumerate}

\textbf{Why SQMS?} The SQMS Center uniquely possesses: (a)~the world's
highest-$Q_0$ nitrogen-doped niobium cavities; (b)~millikelvin cryogenic
infrastructure (dilution refrigerators achieving 10--20~mK);
(c)~decades of expertise in SRF cavity characterisation under Anna
Grassellino's leadership; (d)~existing low-level RF (LLRF) and readout
systems. No other facility in the world can execute this measurement.

\textbf{Why now?} The SQMS Center has been renewed with \$125M in DOE
funding for its second five-year phase (2025--2030), with expanded scope
in quantum sensing. This proposal aligns directly with the SQMS~2.0
mission of ``advancing quantum technology through application to DOE's
most pressing scientific challenge areas.'' The measurement leverages
existing hardware and expertise, requiring only incremental investment in
readout upgrades and dedicated beam time.

\textbf{What we learn either way:}
\begin{itemize}[nosep]
\item \textbf{Positive result} ($Q$-ratio consistent with DFD): First
  laboratory evidence for a new fundamental scalar field. Opens a
  \$100M+ follow-on programme of cavity-based scalar-field physics.
\item \textbf{Null result}: World's best constraint $\lbone <
  7.8\ee{-10}$ (realistic) to $2.5\ee{-10}$ (optimistic).
  Establishes the systematic floor for SRF cavity cross-correlations.
  Directly valuable for SQMS quantum coherence research (identifying
  previously uncharacterised loss mechanisms).
\end{itemize}


%=============================================================================
\section{Scientific Motivation}
\label{sec:motivation}
%=============================================================================

\subsection{The scalar-field hypothesis}

The DFD framework describes gravity through a scalar refractive index
$n(\mathbf{x},t) = e^{\psi(\mathbf{x},t)}$ on flat Minkowski spacetime.
The field equation:
\begin{equation}
\nabla \cdot \left[\mu\!\left(\frac{|\nabla\psi|}{a_*}\right)
\nabla\psi\right] = -\frac{8\pi G}{c^2}\,\rho_{\rm eff},
\label{eq:field_equation}
\end{equation}
with the interpolating function $\mu(x) = x/(1+x)$, recovers Newtonian
gravity at high accelerations and MOND phenomenology at low accelerations.
The theory uses four axioms and one free parameter ($\lambda$, the
EM--scalar coupling constant).

The two-component decomposition $\psi = \psi_{\rm grav} + \delta\psi_{\rm EM}$
separates the gravitational background from the electromagnetically
sourced perturbation. The EM energy density stored in an SRF cavity
sources a local $\delta\psi_{\rm EM}$ that modifies the cavity's
effective surface resistance.

\subsection{Why multi-frequency Q-ratios are decisive}

The DFD scalar field contributes an additional loss channel to SRF
cavities:
\begin{equation}
\frac{1}{Q_\psi(\omega)} = \frac{(\lambda-1)^2\,\Heff^2(\omega)}{\Qpsi}
\cdot\frac{4\pi G\,u_{\rm EM}}{\muG\,c^2},
\label{eq:psi_loss}
\end{equation}
where $\Heff(\omega)$ is the overlap integral between the EM mode and
the psi-mode at frequency $\omega$, $\Qpsi$ is the scalar-field quality
factor, $u_{\rm EM}$ is the stored EM energy density, and $\muG \approx
0.657$ is the galactic-scale MOND interpolation parameter.

The critical insight is that the frequency dependence of the psi-channel
loss enters through $\Heff^2(\omega)$---the geometrical overlap between
the EM mode pattern and the scalar mode pattern. This is
\textit{fundamentally different} from all conventional loss mechanisms:

\begin{table}[H]
\centering
\caption{Frequency scaling of cavity loss mechanisms in a TESLA 9-cell
cavity. The psi-channel signature is dramatically distinct from all
conventional mechanisms.}
\label{tab:freq_scaling}
\renewcommand{\arraystretch}{1.3}
\begin{tabular}{@{}llccc@{}}
\toprule
\textbf{Loss mechanism} & \textbf{Scaling} &
  $\dfrac{Q_0^{-1}(2.4\;\text{GHz})}{Q_0^{-1}(1.3\;\text{GHz})}$ &
  $\dfrac{Q_0^{-1}(1.7\;\text{GHz})}{Q_0^{-1}(1.3\;\text{GHz})}$ &
  \textbf{Ref.} \\
\midrule
BCS surface resistance & $\propto \omega^2$ & 3.41 & 1.71 & Mattis--Bardeen \\
Residual resistance & Constant & 1.00 & 1.00 & Empirical \\
Trapped flux & $\propto \omega$ (approx.) & 1.85 & 1.31 & Gurevich \\
\rowcolor{yellow!20}
\textbf{DFD psi-channel} & $\propto \Heff^2(\omega)$ & \textbf{0.49} & \textbf{0.25} & This proposal \\
\bottomrule
\end{tabular}
\end{table}

\noindent The psi-channel loss \textit{decreases} at higher frequencies
(because higher-order modes have field distributions that are more
concentrated on the cavity axis, reducing the volume overlap with the
scalar mode). This is the opposite of BCS behaviour. A measurement of
the multi-frequency $Q_0$ ratio to $\sim$10\% precision would
unambiguously identify or exclude a psi-channel contribution.

\textbf{This is the single most powerful diagnostic available in the
entire DFD detection programme.}

\subsection{Current constraints and the discovery window}

The current best constraint on the DFD coupling constant is:
\begin{equation}
\lbone < 1.56\ee{-9} \quad\text{(SQMS, single-cavity,
$\mu_0$-corrected)},
\end{equation}
derived from the absence of anomalous residual resistance in
nitrogen-doped niobium cavities. This bound is achieved passively: the
existing SQMS $Q_0$ data already constrain $\lambda$.

The proposed experiment improves this bound through two mechanisms:
\begin{enumerate}[nosep]
\item \textbf{Statistical improvement}: Cross-correlating four cavities
  over 10 science runs yields an effective averaging factor
  $\sqrt{N_{\rm eff}}$, where $N_{\rm eff} = 4$--40 depending on
  systematic control. This improves the bound by a factor of 2--6.3$\times$.
\item \textbf{Frequency diagnostic}: The multi-frequency $Q$-ratio test
  is qualitatively new. It converts a \textit{constraint} (absence of
  anomalous loss) into a \textit{diagnostic} (positive identification of
  the loss mechanism). Even if the bound improvement is modest, the
  frequency ratio measurement is decisive.
\end{enumerate}

The discovery window spans $7.8\ee{-10} < \lbone < 1.56\ee{-9}$
(factor-of-2 improvement, conservative) to
$2.5\ee{-10} < \lbone < 1.56\ee{-9}$ (factor-of-6 improvement,
optimistic). Any detection within this window would be the first
laboratory evidence for a new fundamental scalar field.


%=============================================================================
\section{Theoretical Context: Why This Matters Beyond DFD}
\label{sec:broader_context}
%=============================================================================

The proposed measurement is valuable independent of DFD's ultimate
validity:

\begin{enumerate}
\item \textbf{Generic scalar-field tests.} Any theory with a scalar
  field $\phi$ coupled to the electromagnetic Lagrangian
  $\mathcal{L}_{\rm EM} = -(1/4)f(\phi)F_{\mu\nu}F^{\mu\nu}$
  produces frequency-dependent cavity losses. Our measurement constrains
  the coupling function $f(\phi)$ for \textit{any} such theory, including
  dilaton models, chameleon fields, and axion-like particles with
  photon coupling.

\item \textbf{Precision SRF surface physics.} The multi-frequency
  $Q_0$-ratio measurement is the most precise test of the BCS
  frequency scaling ever performed below 100~mK. Any deviation---from
  DFD or otherwise---represents new surface physics. This directly
  benefits SQMS's core mission of understanding and eliminating
  decoherence sources in superconducting quantum devices.

\item \textbf{Quantum sensing demonstration.} SRF cavities operating at
  $Q_0 > 2\ee{11}$ are among the most sensitive electromagnetic
  resonators ever built. Demonstrating their use as fundamental physics
  sensors extends the SQMS programme into quantum sensing for
  fundamental science, a strategic priority for DOE.

\item \textbf{Pathfinder for dark matter searches.} SRF cavity searches
  for dark photons and axion-like particles (ALPs) use similar
  multi-frequency and cross-correlation techniques. Our systematic
  characterisation directly supports future SQMS dark-sector searches.
\end{enumerate}


%=============================================================================
%=============================================================================
\part{Experimental Design}
%=============================================================================
%=============================================================================

%=============================================================================
\section{Cavity Array Configuration}
\label{sec:cavity_config}
%=============================================================================

\subsection{The $2\times 2$ arrangement}

Four TESLA-type 9-cell niobium cavities are arranged in two pairs:

\begin{itemize}
\item \textbf{Pair~A} (Cavities 1 and 2): Installed in the same
  cryostat, sharing the same vacuum space, with $\sim$0.5~m centre-to-centre
  separation. Both cavities share all environmental systematics:
  temperature fluctuations, magnetic field drifts, mechanical
  vibrations, and He pressure variations.

\item \textbf{Pair~B} (Cavities 3 and 4): Installed in a separate
  cryostat (or the same cryostat but a different vacuum space), with
  $>$2~m separation from Pair~A. Pair~B has \textit{independent}
  environmental systematics from Pair~A.
\end{itemize}

This arrangement enables six cross-correlations with three
qualitatively different characters:

\begin{table}[H]
\centering
\caption{Cross-correlation channels and their diagnostic power.}
\label{tab:xcorr}
\renewcommand{\arraystretch}{1.2}
\begin{tabular}{@{}llll@{}}
\toprule
\textbf{Channel} & \textbf{Type} & \textbf{Psi signal?} &
  \textbf{Environmental?} \\
\midrule
$C_{12}$ & Intra-pair A & Yes (strong) & Yes (correlated) \\
$C_{34}$ & Intra-pair B & Yes (strong) & Yes (correlated) \\
$C_{13}, C_{14}$ & Inter-pair & Yes (strong) & \textbf{No} (independent) \\
$C_{23}, C_{24}$ & Inter-pair & Yes (strong) & \textbf{No} (independent) \\
\bottomrule
\end{tabular}
\end{table}

\noindent\textbf{Key principle}: A genuine psi-field signal produces
equal correlations in all six channels ($C_{12} = C_{34} = C_{13}
= \ldots$), because the scalar field is spatially coherent over the
$\sim$2~m array baseline (wavelength $\gg 2$~m for cosmological
sources). Environmental systematics produce $C_{12}, C_{34} \gg C_{13}$
(intra-pair correlations dominate). This is the primary systematic
rejection mechanism.

\subsection{Cavity specifications}

\begin{table}[H]
\centering
\caption{Required cavity specifications. All parameters are within
demonstrated SQMS capabilities.}
\label{tab:cavity_specs}
\renewcommand{\arraystretch}{1.25}
\begin{tabular}{@{}lll@{}}
\toprule
\textbf{Parameter} & \textbf{Specification} & \textbf{Demonstrated?} \\
\midrule
Material & Bulk Nb, RRR $\geq$ 300 & Yes (standard SQMS) \\
Surface treatment & Nitrogen-doped (800$^\circ$C/20~min N$_2$, EP) &
  Yes (Grassellino protocol) \\
Geometry & TESLA 9-cell, 1.3~GHz fundamental & Yes (SQMS inventory) \\
$Q_0$ at 2~K, 2~MV/m & $\geq 2\ee{10}$ & Yes ($5\ee{10}$ demonstrated) \\
$Q_0$ at 20~mK, 2~MV/m & $\geq 2\ee{11}$ (target) & Demonstrated at SQMS \\
Residual resistance & $R_{\rm res} < 2$~n$\Omega$ & Yes (N-doped) \\
Trapped flux & $B_{\rm trapped} < 1$~nT & Yes ($\mu$-metal shielding) \\
Higher-order modes & TM$_{011}$ ($\sim$2.4~GHz), TE$_{011}$
  ($\sim$1.7~GHz) accessible & Yes \\
\bottomrule
\end{tabular}
\end{table}

\subsection{Cryogenic configuration}

\textbf{Primary operating point: 20~mK in dilution refrigerators.}

At 20~mK, BCS surface resistance becomes negligible:
\begin{equation}
R_s^{\rm BCS}(20\;\text{mK}, 1.3\;\text{GHz}) \approx
A\,\omega^2\,\frac{1}{\kB T}\,
\exp\!\left(-\frac{\Delta_0}{\kB T}\right)
\ll R_{\rm res},
\end{equation}
where $\Delta_0/\kB \approx 17.1$~K for Nb. At 20~mK, the
BCS contribution is suppressed by $\exp(-17.1/0.02) = \exp(-855) \approx 0$.
The measured $Q_0$ is therefore dominated by the residual resistance
$R_{\rm res}$ and any anomalous (psi-channel) contribution.

This is optimal for our measurement: the BCS ``pedestal'' is removed,
making the residual and any psi-channel loss the dominant physics. Any
frequency-dependent contribution to the loss that does \textit{not}
follow the BCS $\omega^2$ scaling is immediately suspicious.

\textbf{Secondary operating point: 1.8~K (liquid helium bath).}

At 1.8~K, BCS resistance is significant ($R_s^{\rm BCS} \sim
10$~n$\Omega$ at 1.3~GHz). The BCS contribution serves as a
\textit{calibration standard}: the known $\omega^2$ scaling provides an
in-situ verification of the frequency-dependent measurement systematics.
If our multi-frequency readout correctly reproduces $\omega^2$ scaling
at 1.8~K (where BCS dominates), we have high confidence in anomalous
results at 20~mK (where BCS is absent).

\textbf{Temperature scan protocol:} Measurements at 8 temperature
points (20, 50, 100, 200, 500~mK; 1.2, 1.5, 1.8~K) map the full
$R_s(T)$ curve and verify the BCS gap parameter $\Delta_0$.


%=============================================================================
\section{Multi-Frequency Measurement Strategy}
\label{sec:multi_freq}
%=============================================================================

\subsection{Mode selection and overlap integrals}

Each TESLA 9-cell cavity supports multiple electromagnetic modes.
We target three modes that provide maximum diagnostic power:

\begin{table}[H]
\centering
\caption{Target electromagnetic modes and their psi-field overlap
integrals $\Heff$.}
\label{tab:modes}
\renewcommand{\arraystretch}{1.25}
\begin{tabular}{@{}lcccc@{}}
\toprule
\textbf{Mode} & \textbf{Frequency} & \textbf{Field pattern} &
  $\Heff$ & $\Heff^2/\Heff^2(\text{TM}_{010})$ \\
\midrule
TM$_{010}$ & 1.3~GHz & Max $E_z$ on axis & $3.0\ee{-5}$ & 1.00 \\
TE$_{011}$ & $\sim$1.7~GHz & Max $E$ at equator & $1.5\ee{-5}$ & 0.25 \\
TM$_{011}$ & $\sim$2.4~GHz & Concentrated on axis & $2.1\ee{-5}$ & 0.49 \\
\bottomrule
\end{tabular}
\end{table}

\noindent The overlap integral $\Heff(\omega) = \int_V |\mathbf{E}_{\rm EM}
(\mathbf{r};\omega)|^2\,u_\psi(\mathbf{r};\omega)\,d^3r / \sqrt{\int
|\mathbf{E}|^4\,d^3r\,\int u_\psi^2\,d^3r}$ is computed from finite-element
simulations of the TESLA cavity geometry. The psi-mode profile $u_\psi$
is determined by the scalar field equation in the cavity geometry; for
the fundamental scalar mode, $u_\psi$ is approximately uniform over the
cavity volume, and the overlap is dominated by the EM mode's volume
integral $\int |\mathbf{E}|^2\,d^3r$.

\subsection{The decisive Q-ratio diagnostic}

\begin{tcolorbox}[colback=yellow!10!white, colframe=sqmsgold,
  title=\textbf{THE KEY MEASUREMENT}]
Measure the ratio of excess (non-BCS, non-residual) losses at two
frequencies:
\begin{equation}
\boxed{
\mathcal{R} \equiv \frac{\delta Q_0^{-1}(2.4~\text{GHz})}
{\delta Q_0^{-1}(1.3~\text{GHz})}
}
\end{equation}
where $\delta Q_0^{-1} \equiv Q_0^{-1}(\text{measured}) -
Q_0^{-1}(\text{BCS}) - R_{\rm res}/G_{\rm geom}$ is the excess loss
after subtracting known BCS and residual contributions.

\textbf{DFD prediction}: $\mathcal{R}_{\rm DFD} = \Heff^2(2.4)/
\Heff^2(1.3) = 0.49 \pm 0.05$ (geometry-dependent uncertainty).

\textbf{BCS prediction} (if residual resistance has $\omega^2$ component):
$\mathcal{R}_{\rm BCS} = (2.4/1.3)^2 = 3.41$.

\textbf{Frequency-independent residual}: $\mathcal{R}_{\rm res} = 1.00$.

\textbf{Trapped flux}: $\mathcal{R}_{\rm flux} \approx 1.85$.

These four predictions are separated by factors of 2--7$\times$. A
10\% measurement of $\mathcal{R}$ decisively distinguishes all four
hypotheses.
\end{tcolorbox}

\subsection{Measurement of Q$_0$ at each frequency}

For each mode, $Q_0$ is determined by the ring-down method:
\begin{enumerate}[nosep]
\item Drive the cavity to the target stored energy $U_0$ at gradient
  $E_{\rm acc} = 2$~MV/m.
\item Switch off the drive and measure the exponential decay of transmitted
  power: $P(t) = P_0\,e^{-2\pi f_0 t/Q_0}$.
\item Extract $Q_0$ from the time constant $\tau = Q_0/(\pi f_0)$.
\item Repeat for $N \geq 100$ ring-downs per measurement point.
\end{enumerate}

The per-ring-down precision is:
\begin{equation}
\frac{\sigma_Q}{Q_0}\bigg|_{\rm single} = \frac{1}{\sqrt{2\pi f_0\,\tau}}
\cdot\frac{1}{\SNR_{\rm readout}}.
\end{equation}
For TM$_{010}$ at 1.3~GHz with $Q_0 = 2\ee{11}$: $\tau = 49$~s,
$\SNR_{\rm readout} = 10^3$ (achievable with existing LLRF systems),
giving $\sigma_Q/Q_0|_{\rm single} \sim 1.6\ee{-9}$. Over 100
ring-downs: $\sigma_Q/Q_0 \sim 1.6\ee{-10}$.

For TM$_{011}$ at 2.4~GHz with $Q_0 \sim 5\ee{10}$ (lower due to
$\omega^2$ BCS at non-zero $T$): $\tau \sim 6.6$~s, giving
$\sigma_Q/Q_0|_{\rm single} \sim 3\ee{-9}$. Over 100 ring-downs:
$\sigma_Q/Q_0 \sim 3\ee{-10}$.


%=============================================================================
\section{RF Electronics and Data Acquisition}
\label{sec:electronics}
%=============================================================================

\subsection{Phase-locked RF system}

Each cavity requires an independent RF drive and readout chain:

\begin{enumerate}[nosep]
\item \textbf{Master oscillator}: Ultra-low-phase-noise synthesiser at
  1.3~GHz, distributed to all four cavities. Phase stability: $<1$~fs
  jitter over 1000~s (achievable with commercial Wenzel or Pascall
  oscillators).

\item \textbf{Frequency multiplexing}: For multi-frequency operation,
  each cavity is driven sequentially at TM$_{010}$, TE$_{011}$, and
  TM$_{011}$. Mode switching time: $<1$~s (retuning the drive frequency
  and re-locking the phase-locked loop). The three modes are
  non-degenerate and well-separated in frequency ($>$400~MHz spacing),
  ensuring clean mode identification.

\item \textbf{Pickup antenna}: Weakly coupled antenna ($\beta < 10^{-3}$)
  at each cavity, feeding a cryogenic low-noise amplifier (LNA) at
  4~K. Noise temperature $T_N < 5$~K, ensuring readout SNR $> 10^3$
  at the target stored energy.

\item \textbf{Digitiser}: 16-bit ADC at 10~MS/s per channel, with
  GPS-disciplined clock for absolute timing. Four simultaneous channels
  (one per cavity). Continuous data streaming to local storage.
\end{enumerate}

\subsection{Data acquisition and storage}

\begin{itemize}[nosep]
\item \textbf{Raw data rate}: 4 channels $\times$ 10~MS/s $\times$
  2~bytes = 80~MB/s = 6.9~TB/day. Stored for 30 days of science runs
  $\times$ 10 runs = 2.1~PB total. \textit{Note}: Ring-down analysis
  reduces this to $\sim$10~GB of processed $Q_0(t)$ time series.
\item \textbf{Environmental data}: 20 sensor channels (temperature,
  magnetic field, vibration, pressure) at 1~kHz sampling = 144~GB/day.
\item \textbf{Processing}: Real-time ring-down fitting on a
  dedicated GPU node. Offline cross-correlation analysis on Fermilab
  computing cluster.
\end{itemize}


%=============================================================================
%=============================================================================
\part{Systematic Error Budget}
%=============================================================================
%=============================================================================

%=============================================================================
\section{Identified Systematics and Mitigation}
\label{sec:systematics}
%=============================================================================

\begin{longtable}{@{}p{2.8cm}p{2.5cm}p{2cm}p{2.5cm}p{3cm}@{}}
\caption{Complete systematic error budget for SQMS Phase~I. Each
systematic is classified by magnitude, frequency character, and
distinguishability from a psi signal.}
\label{tab:systematics}\\
\toprule
\textbf{Systematic} & \textbf{Magnitude} ($\delta Q/Q$) &
  \textbf{Freq.\ dep.} & \textbf{Psi-like?} &
  \textbf{Mitigation} \\
\midrule
\endfirsthead
\toprule
\textbf{Systematic} & \textbf{Magnitude} ($\delta Q/Q$) &
  \textbf{Freq.\ dep.} & \textbf{Psi-like?} &
  \textbf{Mitigation} \\
\midrule
\endhead
\bottomrule
\endlastfoot
%
Temperature drift &
  $\sim 10^{-4}$ per $\mu$K at 20~mK &
  Follows BCS ($\omega^2$) &
  \textbf{No} (wrong freq.\ dep.) &
  $\delta T/T < 10^{-4}$; regression on $T(t)$; inter-pair decorrelation \\
\midrule
Trapped magnetic flux &
  $\sim 5\ee{-4}$ per nT &
  $\propto \omega$ (approx.) &
  Partially &
  $B_{\rm res} < 1$~nT; continuous 3-axis fluxgate; field compensation coils \\
\midrule
Lorentz detuning &
  $\delta f/f \sim 10^{-12}$ at 2~MV/m &
  $\propto E_{\rm acc}^2$ &
  No &
  Operate at low gradient (2~MV/m); piezo tuner feedback \\
\midrule
Mechanical vibration &
  $\delta Q/Q \sim 10^{-8}$ per nm at 1~Hz &
  Broadband &
  No &
  Passive + active isolation ($<10^{-9}$ m/$\sqrt{\text{Hz}}$);
  microphonics detuning monitor \\
\midrule
Cosmic ray bursts &
  $\sim 10^{-3}$ per event, $\sim$0.1~Hz rate &
  Impulsive &
  \textbf{Partially} (correlated in Pair~A) &
  Scintillator veto; time-domain template subtraction; decorrelated between
  Pairs A and B \\
\midrule
Helium pressure fluctuations &
  $\delta Q/Q \sim 10^{-7}$ per $\mu$bar &
  Low-$f$ drift &
  No &
  Pressure stabilisation; regression on $P(t)$ \\
\midrule
Readout chain systematics &
  $\delta Q/Q \sim 10^{-9}$ (calibration) &
  Systematic offset &
  No &
  Cross-calibration using known BCS $\omega^2$ scaling at 1.8~K \\
\midrule
Surface chemistry evolution &
  Slow drift over weeks &
  Unknown &
  Potentially &
  Interleave baseline re-measurements; compare Pairs A and B drift
  rates \\
\bottomrule
\end{longtable}

\subsection{The dominant systematic: temperature drift}

At 20~mK, the BCS resistance is exponentially suppressed, but the
\textit{derivative} $\partial R_s/\partial T$ is also suppressed
(the exponential kills both the function and its derivative). The
dominant temperature-induced $Q_0$ drift at 20~mK comes through the
\textit{residual resistance} temperature coefficient:
\begin{equation}
\frac{\partial R_{\rm res}}{\partial T} \sim 0.01~\text{n}\Omega/\text{mK}
\quad\text{(empirical, N-doped Nb)}.
\end{equation}
For $\delta T = 1~\mu$K: $\delta R_{\rm res} \sim 10^{-5}$~n$\Omega$,
giving $\delta Q/Q \sim 10^{-5}\times Q_0/(G_{\rm geom}/R_{\rm res})
\sim 10^{-8}$.

This is well below the statistical precision ($10^{-10}$ per
measurement). The 2$\times$2 arrangement provides additional rejection:
temperature-induced correlations appear in $C_{12}$ (same cryostat) but
not in $C_{13}$ (different cryostats), enabling model-independent
subtraction.

\subsection{The secondary systematic: trapped magnetic flux}

Trapped magnetic flux produces a frequency-dependent loss
$R_{\rm flux} \propto B_{\rm trapped}\times\omega$ (approximate, with
detailed frequency dependence depending on the pinning landscape).
This has a different frequency scaling from both BCS ($\omega^2$) and
psi-channel ($\Heff^2$), but is closer to the psi-channel signature
than BCS.

\textbf{Mitigation}:
\begin{enumerate}[nosep]
\item $\mu$-metal shielding to $B_{\rm res} < 1$~nT (SQMS standard).
\item Active field compensation with Helmholtz coils, monitored by
  3-axis fluxgate magnetometers with 10~pT resolution.
\item The \textit{ratio} $\Heff^2(2.4)/\Heff^2(1.3) = 0.49$ is
  distinct from the trapped-flux ratio $\sim (2.4/1.3) = 1.85$.
  A 20\% measurement of $\mathcal{R}$ distinguishes them.
\end{enumerate}


%=============================================================================
%=============================================================================
\part{Detection Criteria and Analysis}
%=============================================================================
%=============================================================================

%=============================================================================
\section{Four-Criterion Detection Threshold}
\label{sec:detection_criteria}
%=============================================================================

A DFD psi-field detection is claimed if and only if \textbf{all four}
of the following criteria are satisfied simultaneously:

\begin{tcolorbox}[colback=green!5!white, colframe=green!60!black,
  title=\textbf{DETECTION CRITERIA (5$\sigma$ discovery)}]
\begin{enumerate}[label=\textbf{C\arabic*.}, itemsep=8pt]

\item \textbf{Inter-pair cross-correlation significance.}
  The environmentally cleaned inter-pair cross-correlation is non-zero:
  \begin{equation}
  \tilde{C}_{13}(0) > 5\,\sigma_{\tilde{C}},
  \end{equation}
  where $\tilde{C}_{ij}$ is computed after regression on all monitored
  environmental channels ($T, B, x, P$).

\item \textbf{Intra/inter-pair consistency.}
  The intra-pair and inter-pair correlations are statistically consistent:
  \begin{equation}
  |\tilde{C}_{12}(0) - \tilde{C}_{13}(0)| < 2\,\sigma_{\rm diff},
  \end{equation}
  ruling out systematics that correlate cavities within a cryostat
  but not between cryostats.

\item \textbf{Multi-frequency Q-ratio.}
  The measured excess-loss ratio $\mathcal{R}$ is consistent with the
  DFD prediction and inconsistent with BCS:
  \begin{equation}
  |\mathcal{R}_{\rm meas} - 0.49| < 3\,\sigma_{\mathcal{R}}
  \quad\text{AND}\quad
  |\mathcal{R}_{\rm meas} - 3.41| > 5\,\sigma_{\mathcal{R}}.
  \end{equation}

\item \textbf{Detuning null test.}
  When one cavity of a pair is detuned by $>$10 bandwidths from
  resonance (while the other remains on resonance), the cross-correlation
  vanishes:
  $\tilde{C}_{ij}^{\rm detuned}(0)$ consistent with zero at $2\sigma$.
  This exploits the fact that psi-coupling is resonant (the scalar field
  mode must match the EM mode frequency), while environmental couplings
  are broadband.

\end{enumerate}
\end{tcolorbox}

\subsection{Null result interpretation}

If any of the four criteria fails after the complete science campaign:

\begin{equation}
\boxed{
\lbone < \frac{1.56\ee{-9}}{\sqrt{N_{\rm eff}}}
}
\end{equation}

where $N_{\rm eff}$ is the effective number of independent measurements:

\begin{table}[H]
\centering
\caption{Null result bounds by systematic scenario.}
\label{tab:null_bounds}
\renewcommand{\arraystretch}{1.2}
\begin{tabular}{@{}lcc@{}}
\toprule
\textbf{Scenario} & $N_{\rm eff}$ & \textbf{Bound $\lbone$} \\
\midrule
Conservative (strong systematics) & 2 & $1.1\ee{-9}$ \\
Realistic (moderate systematics) & 4 & $7.8\ee{-10}$ \\
Optimistic (statistics-limited) & 40 & $2.5\ee{-10}$ \\
\bottomrule
\end{tabular}
\end{table}

\textbf{Even the conservative null result improves the SQMS bound by
$\sim$40\%.}

\subsection{Blind analysis protocol}

To prevent analyst bias, we adopt the following blinding procedure:
\begin{enumerate}[nosep]
\item The $Q_0(t)$ time series for each cavity are independently
  processed by two analysis teams.
\item A random, secret offset is added to each cavity's $Q_0$ values
  before cross-correlation. The offset is removed only after the
  analysis cuts are frozen.
\item The multi-frequency ratio $\mathcal{R}$ is computed separately by
  both teams before unblinding.
\item Results are unblinded at a single event, with pre-registered
  detection criteria (C1--C4 above).
\end{enumerate}


%=============================================================================
%=============================================================================
\part{Project Plan}
%=============================================================================
%=============================================================================

%=============================================================================
\section{Timeline: 24-Month Programme}
\label{sec:timeline}
%=============================================================================

\begin{table}[H]
\centering
\caption{Phase~I timeline. Each phase includes built-in contingency.}
\label{tab:timeline}
\renewcommand{\arraystretch}{1.3}
\begin{tabular}{@{}clp{8cm}@{}}
\toprule
\textbf{Months} & \textbf{Phase} & \textbf{Activities} \\
\midrule
1--3 & \textbf{Commissioning} &
  Select 4 cavities from SQMS inventory (best $Q_0$, matched RRR).
  Install in $2\times 2$ configuration. Upgrade RF readout for
  multi-frequency operation. Install environmental monitoring
  (fluxgates, accelerometers, thermometry). Commission DAQ system. \\
\midrule
4--6 & \textbf{Baseline} &
  Phase~I-A: Individual cavity $Q_0$ mapping at 5 gradients, 8
  temperatures, 5 magnetic field values ($8\times 5\times 5 = 200$
  measurement points per cavity, 800 total). Establish BCS model
  parameters ($\Delta_0, R_{\rm res}$) for each cavity at each mode.
  Measure environmental transfer functions. \\
\midrule
7--9 & \textbf{Multi-freq.\ calibration} &
  Multi-frequency $Q_0$ ratio measurements at 1.8~K (BCS-dominated).
  Verify $\omega^2$ scaling to $<$5\%. This is the \textbf{systematics
  calibration phase}: if the ratio at 1.8~K deviates from BCS
  prediction, the systematic must be identified before proceeding. \\
\midrule
10--18 & \textbf{Science data} &
  Phase~I-B: 10 science runs of $\geq 10^6$~s ($\sim$12 days) each at
  20~mK, $E_{\rm acc} = 2$~MV/m, $B < 1$~nT. Continuous $Q_0(t)$
  monitoring at all three frequencies. Interleaved detuning null tests
  (Criterion C4). Environmental data recording. \\
\midrule
19--21 & \textbf{Analysis} &
  Blind cross-correlation analysis. Environmental regression.
  Multi-frequency ratio computation. Independent verification by
  second analysis team. Pre-registered criteria evaluation. \\
\midrule
22--24 & \textbf{Publication} &
  Unblinding. Internal review. Publication of results (detection or
  null constraint). Phase~II/III planning based on results. \\
\bottomrule
\end{tabular}
\end{table}


%=============================================================================
\section{Budget}
\label{sec:budget}
%=============================================================================

\begin{table}[H]
\centering
\caption{Detailed budget breakdown. All costs in FY2027 dollars.}
\label{tab:budget}
\renewcommand{\arraystretch}{1.25}
\begin{tabular}{@{}p{5cm}rrl@{}}
\toprule
\textbf{Item} & \textbf{Cost (\$K)} & \textbf{Category} &
  \textbf{Notes} \\
\midrule
\multicolumn{4}{@{}l}{\textbf{A. Personnel (2 years)}} \\
\quad PI (50\% FTE $\times$ 2 yr) & 250 & Labour &
  Experiment design, oversight \\
\quad Co-PI (50\% FTE $\times$ 2 yr) & 250 & Labour &
  Theory interface, analysis \\
\quad Postdoc \#1 (100\% FTE $\times$ 2 yr) & 220 & Labour &
  RF systems, cavity operations \\
\quad Postdoc \#2 (100\% FTE $\times$ 2 yr) & 220 & Labour &
  Data analysis, cross-correlations \\
\quad Technician (100\% FTE $\times$ 2 yr) & 180 & Labour &
  Cryogenics, installation \\
\quad Graduate student \#1 (2 yr) & 70 & Labour &
  Multi-frequency readout development \\
\quad Graduate student \#2 (2 yr) & 70 & Labour &
  Environmental monitoring, systematics \\
\quad Fringe benefits (30\%) & 378 & Labour & \\
\midrule
\textit{Personnel subtotal} & \textit{1,638} & & \\
\midrule
\multicolumn{4}{@{}l}{\textbf{B. Equipment}} \\
\quad Multi-frequency LLRF upgrade (4 channels) & 300 & Equipment &
  Phase-locked loops, synthesisers \\
\quad Cryogenic LNAs (4$\times$ at 4~K) & 80 & Equipment &
  HEMT amplifiers, 1--3~GHz \\
\quad 16-bit ADC system (4 channels) & 60 & Equipment &
  GPS-disciplined, 10~MS/s \\
\quad 3-axis fluxgate magnetometers (4$\times$) & 40 & Equipment &
  10~pT resolution \\
\quad Cryogenic accelerometers (4$\times$) & 30 & Equipment &
  $<10^{-10}$~g/$\sqrt{\text{Hz}}$ \\
\quad Helmholtz compensation coils & 20 & Equipment &
  Active field control \\
\quad $\mu$-metal shielding (upgrade) & 60 & Equipment &
  $B_{\rm res} < 0.5$~nT \\
\quad High-precision thermometry & 25 & Equipment &
  $<1~\mu$K at 20~mK \\
\quad Scintillator cosmic ray veto & 15 & Equipment &
  Coincidence trigger \\
\quad Cables, connectors, RF components & 30 & Equipment & \\
\quad Computing hardware (GPU node + storage) & 40 & Equipment &
  Real-time ring-down fitting \\
\midrule
\textit{Equipment subtotal} & \textit{700} & & \\
\midrule
\multicolumn{4}{@{}l}{\textbf{C. Cryogenics and Operations}} \\
\quad DR operating cost (2 yr) & 200 & Operations &
  Helium, power, maintenance \\
\quad LHe supply for 1.8~K runs & 40 & Operations &
  Baseline calibration \\
\quad Cavity handling and preparation & 30 & Operations &
  Clean-room time \\
\midrule
\textit{Operations subtotal} & \textit{270} & & \\
\midrule
\multicolumn{4}{@{}l}{\textbf{D. Travel and Publication}} \\
\quad Conference travel (6 trips $\times$ 2 yr) & 48 & Travel &
  SRF, QIS, APS meetings \\
\quad Collaboration meetings & 20 & Travel & \\
\quad Publication charges & 10 & Other & \\
\midrule
\textit{Travel/publication subtotal} & \textit{78} & & \\
\midrule
\multicolumn{4}{@{}l}{\textbf{E. Indirect Costs}} \\
\quad Fermilab overhead (20\% on direct) & 537 & Indirect & \\
\midrule
\textbf{TOTAL} & \textbf{\$3,223K} & & \\
\bottomrule
\end{tabular}
\end{table}

\noindent\textbf{Cost efficiency}: The total cost of \$3.2M is modest for
a fundamental physics experiment because:
\begin{enumerate}[nosep]
\item The four cavities already exist in SQMS inventory (\$0 cavity cost).
\item The dilution refrigerators already exist at Fermilab (\$0 cryogenic
  capital).
\item The LLRF expertise already exists in the SQMS team.
\item The primary cost is \textit{people}, not hardware.
\end{enumerate}


%=============================================================================
\section{Personnel and Expertise}
\label{sec:personnel}
%=============================================================================

\subsection{Required expertise}

\begin{table}[H]
\centering
\caption{Required expertise and potential SQMS contributors.}
\label{tab:personnel}
\renewcommand{\arraystretch}{1.2}
\begin{tabular}{@{}p{3cm}p{4cm}p{5.5cm}@{}}
\toprule
\textbf{Role} & \textbf{Expertise needed} & \textbf{SQMS capability} \\
\midrule
PI & SRF cavity physics; precision $Q_0$ measurements; cryogenics
  leadership &
  SQMS has world-leading SRF expertise under Grassellino. Multiple senior
  scientists qualified (Romanenko, Bafia, Murthy). \\
\midrule
Co-PI (Theory) & Scalar-field theory; DFD formalism; statistical analysis
  of cross-correlation data &
  External collaborator. Interface to DFD theoretical predictions. \\
\midrule
Postdoc~1 & RF engineering; phase-locked loops; multi-frequency cavity
  characterisation &
  SQMS regularly hires RF postdocs. LLRF group at Fermilab. \\
\midrule
Postdoc~2 & Data science; time-series analysis; blind analysis protocols &
  SQMS data science group. Cross-pollination with LIGO analysis methods. \\
\midrule
Technician & Cryogenic operations; cavity installation; clean-room
  procedures &
  Fermilab SRF infrastructure group. \\
\midrule
Grad students & Multi-frequency readout; environmental monitoring &
  University partners (Northwestern, U.\ Chicago, IIT). \\
\bottomrule
\end{tabular}
\end{table}

\subsection{Institutional resources}

\textbf{SQMS Center at Fermilab} provides:
\begin{itemize}[nosep]
\item Vertical test stands for cavity characterisation
\item Dilution refrigerators capable of $<$10~mK base temperature
\item Clean-room facilities (ISO Class~100) for cavity handling
\item LLRF and controls infrastructure
\item High-purity niobium processing (EP, HPR, N-doping)
\item Magnetic shielding and vibration isolation systems
\item Computing cluster for data analysis
\end{itemize}

\textbf{Consortium partners} (Northwestern, U.\ Chicago, SLAC, NIST)
provide complementary expertise in quantum sensing, precision
metrology, and theoretical physics.


%=============================================================================
%=============================================================================
\part{Strategic Context}
%=============================================================================
%=============================================================================

%=============================================================================
\section{Why This Experiment Over Other DFD Detection Channels}
\label{sec:comparison}
%=============================================================================

\begin{table}[H]
\centering
\caption{Comparison of the top three DFD detection experiments.
SQMS Phase~I is the optimal first experiment on every metric except
raw sensitivity.}
\label{tab:comparison}
\renewcommand{\arraystretch}{1.3}
\begin{tabular}{@{}p{3cm}ccccc@{}}
\toprule
\textbf{Criterion} & \textbf{SQMS Phase~I} & \textbf{LIGO archival} &
  \textbf{Channel~B} \\
\midrule
$\lbone$ reach & $7.8\ee{-10}$ & $\sim 3\ee{-14}$ &
  $\sim 8\ee{-13}$ \\
Cost & \$3.2M & \$0.1--0.5M & \$2--5M \\
Timeline to result & 24 months & 12--18 months & 30--42 months \\
TRL & 8 & 9 & 5 \\
Frequency diagnostic? & \textbf{Yes} (decisive) & No & No \\
Improves current bound? & \textbf{Yes} & \textbf{Yes} (different observable) &
  Depends on $\Qpsi$ \\
Infrastructure exists? & \textbf{Yes} (SQMS) & \textbf{Yes} (LIGO O3 data) &
  Partially \\
Risk & Low & Medium (systematics) & \textbf{High} ($\Qpsi$ unknown) \\
Gateway experiment? & \textbf{Yes} & No & Conditional \\
\bottomrule
\end{tabular}
\end{table}

\textbf{Why SQMS Phase~I first, not LIGO archival?}

The LIGO archival analysis (Rank~\#2) has deeper raw sensitivity but
suffers from two fundamental limitations:
\begin{enumerate}[nosep]
\item It provides a \textit{constraint}, not a \textit{diagnostic}.
  A null result tells you nothing about why the signal is absent.
  A positive result is ambiguous (Newtonian noise, Schumann resonances).
\item It requires a source model (the bound depends on what sourced the
  scalar perturbation). SQMS Phase~I is \textit{source-independent}: it
  measures a cavity property, not an astrophysical signal.
\end{enumerate}

The SQMS multi-frequency Q-ratio is the \textit{only} experiment that
provides a \textbf{positive identification} of the scalar-field loss
mechanism. All other experiments provide bounds. This makes SQMS
Phase~I the indispensable \textbf{gateway experiment}: a positive
Q-ratio result triggers the entire downstream programme (Channel~B,
Phase~II MEMS, Phase~III arrays), while a null result either
(a)~tightens the bound or (b)~identifies the systematic floor for all
future SRF-based searches.

\textbf{Why not Channel~B first?}

Channel~B (SRF source + MEMS detector) was promoted to the
highest-priority bound-improving experiment in the W3i8 definitive
ranking, with $\lbone \sim 8\ee{-13}$. However, the W3i9 audit
revealed a critical dependence on the unknown scalar-field quality
factor $\Qpsi$:
\begin{itemize}[nosep]
\item The $8\ee{-13}$ figure assumes $\Qpsi = 10^{15}$ (unconstrained).
\item Pessimistic reach ($\Qpsi = 10^9$): $\lbone \sim 3\ee{-3}$
  --- cannot improve the current bound.
\item Channel~B should be funded \textit{only after SQMS Phase~I
  provides evidence that the psi-channel exists}.
\end{itemize}


%=============================================================================
\section{Decision Tree: What Happens After Phase~I?}
\label{sec:decision_tree}
%=============================================================================

\begin{tcolorbox}[colback=blue!5!white, colframe=sqmsblue,
  title=\textbf{POST-PHASE~I DECISION TREE}]

\textbf{Branch A: Positive result} (all four criteria met)
\begin{enumerate}[nosep]
\item Extract $(\lambda-1)^2/\Qpsi$ from cross-correlation amplitude.
\item Extract $\Heff^2(\omega)$ ratios from multi-frequency data.
\item $\Qpsi$ correlation timescale provides first empirical estimate of
  $\Qpsi$.
\item \textbf{Immediate follow-on}: Channel~B (\$2--5M) to achieve
  $\lbone \sim 10^{-13}$ using the measured $\Qpsi$.
\item \textbf{Parallel}: LIGO archival analysis with DFD source model
  calibrated by SQMS result.
\item \textbf{Phase~II}: MEMS hybrid (\$3--6M, $\lbone \sim 10^{-15}$).
\item \textbf{Phase~III}: ng-MEMS array (\$5--10M, $\lbone \sim 10^{-15}$)
  or 256-cavity (\$80--120M, $\lbone \sim 10^{-11}$).
\end{enumerate}

\vspace{6pt}
\textbf{Branch B: Null result} (any criterion fails)
\begin{enumerate}[nosep]
\item Publish world-best constraint $\lbone < 7.8\ee{-10}$
  (to $2.5\ee{-10}$).
\item Publish systematic floor characterisation for SRF
  cross-correlations.
\item \textbf{Conditional}: If Q-ratio measurement shows anomalous
  frequency dependence (not BCS, not psi, but unknown), publish as
  new surface physics discovery. This is independently valuable for
  SQMS quantum coherence programme.
\item \textbf{Phase~0 archival}: LIGO archival, multi-GNSS timing, and
  CLONETS-DS proceed on independent track (independent of SQMS result).
\item \textbf{Conditional Phase~II}: ng-MEMS programme (\$5--10M) as
  the cost-optimal path to $\lbone \sim 10^{-15}$, bypassing the
  SRF-based approach.
\end{enumerate}
\end{tcolorbox}


%=============================================================================
\section{Broader Impacts}
\label{sec:broader}
%=============================================================================

\begin{enumerate}
\item \textbf{Workforce development}: Two postdocs and two graduate
  students trained in precision SRF measurements, cryogenic engineering,
  and advanced data analysis. These skills are directly transferable to
  quantum computing (SQMS core mission), accelerator physics, and
  precision measurement science.

\item \textbf{Facility infrastructure}: Multi-frequency readout upgrades
  and environmental monitoring systems installed for this experiment
  become permanent SQMS infrastructure, benefiting all future cavity
  measurements.

\item \textbf{Cross-disciplinary impact}: The multi-frequency Q-ratio
  technique is applicable to dark photon searches, axion-like particle
  searches, and fifth-force constraints using SRF cavities. A single
  infrastructure investment serves multiple fundamental physics
  programmes.

\item \textbf{International collaboration}: DFD theory development
  (Alcock, UK/NZ) provides an international partnership component,
  broadening SQMS's collaboration network.
\end{enumerate}


%=============================================================================
%=============================================================================
\part{Summary and Conclusion}
%=============================================================================
%=============================================================================

This proposal requests \$3.2M over 24 months to execute the single most
cost-effective test of scalar-field modifications to electrodynamics
using the world's most sensitive electromagnetic resonators at the SQMS
Center. The experiment:

\begin{enumerate}
\item \textbf{Leverages existing SQMS assets}: Four N-doped Nb TESLA
  cavities, dilution refrigerators, LLRF systems, and decades of SRF
  expertise. No new cryogenic capital required.

\item \textbf{Deploys a decisive diagnostic}: The multi-frequency Q-ratio
  ($\mathcal{R}_{\rm DFD} = 0.49$ vs $\mathcal{R}_{\rm BCS} = 3.41$)
  is the single most powerful test in the DFD detection programme,
  providing positive identification (not just constraint) of scalar-field
  coupling.

\item \textbf{Guarantees scientific output}: A null result establishes
  $\lbone < 7.8\ee{-10}$ (world best from SRF) and characterises the
  systematic floor for all future cavity-based searches. A positive
  result would be the first laboratory evidence for a new fundamental
  field.

\item \textbf{Opens a programme}: A positive Phase~I result triggers
  Channel~B (\$2--5M), Phase~II MEMS (\$3--6M), and ultimately
  Phase~III arrays (\$5--120M) --- a multi-decade programme of
  increasing precision.

\item \textbf{Aligns with SQMS~2.0}: Quantum sensing for fundamental
  science is a DOE strategic priority. This experiment demonstrates
  SRF cavities as instruments for fundamental physics discovery, beyond
  their traditional role in accelerator science.
\end{enumerate}

\vspace{12pt}
\begin{center}
\fbox{\parbox{0.85\textwidth}{
\centering
\textbf{Total request}: \$3.2M\\
\textbf{Duration}: 24 months\\
\textbf{Key measurement}: Multi-frequency Q-ratio
$\mathcal{R} = 0.49$ (DFD) vs $3.41$ (BCS)\\
\textbf{Sensitivity}: $\lbone < 7.8\ee{-10}$ (null) or detection
(positive)\\
\textbf{Gateway to}: \$100M+ programme of scalar-field physics
}}
\end{center}


%=============================================================================
% REFERENCES
%=============================================================================

\newpage
\begin{thebibliography}{99}

\bibitem{Alcock2026}
G.T.~Alcock, \textit{Density Field Dynamics: A Complete Unified Theory},
v3.2 (2026).

\bibitem{Grassellino2013}
A.~Grassellino et al.,
``Nitrogen and argon doping of niobium for superconducting radio
frequency cavities: a pathway to highly efficient accelerating structures,''
\textit{Supercond.\ Sci.\ Technol.}\ \textbf{26}, 102001 (2013).

\bibitem{Grassellino2017}
A.~Grassellino et al.,
``Unprecedented quality factors at accelerating gradients up to 45 MV/m
in niobium superconducting cavities via low temperature nitrogen infusion,''
\textit{Supercond.\ Sci.\ Technol.}\ \textbf{30}, 094004 (2017).

\bibitem{Romanenko2020}
A.~Romanenko et al.,
``Three-dimensional superconducting resonators at $T < 20$~mK with
photon lifetimes up to $\tau = 2$~s,''
\textit{Phys.\ Rev.\ Applied}\ \textbf{13}, 034032 (2020).

\bibitem{Gurevich2012}
A.~Gurevich,
``Reduction of dissipative nonlinear conductivity of superconductors by
static magnetic field,''
\textit{Phys.\ Rev.\ Lett.}\ \textbf{113}, 087001 (2014).

\bibitem{SQMS2025}
SQMS Center, Fermilab,
``SQMS Phase~2: Shaping the Future of Quantum Information Science,''
DOE Award (2025).

\bibitem{DOE_QIS_2025}
U.S.\ Department of Energy,
``\$625 Million to Advance the Next Phase of National Quantum Information
Science Research Centers'' (November 2025).

\bibitem{Mattis1958}
D.C.~Mattis and J.~Bardeen,
``Theory of the anomalous skin effect in normal and superconducting metals,''
\textit{Phys.\ Rev.}\ \textbf{111}, 412 (1958).

\bibitem{Padamsee2008}
H.~Padamsee, J.~Knobloch, and T.~Hays,
\textit{RF Superconductivity for Accelerators}, 2nd ed.\
(Wiley-VCH, 2008).

\bibitem{McGaugh2016}
S.S.~McGaugh, F.~Lelli, and J.M.~Schombert,
``Radial acceleration relation in rotationally supported galaxies,''
\textit{Phys.\ Rev.\ Lett.}\ \textbf{117}, 201101 (2016).

\bibitem{Bekenstein1984}
J.~Bekenstein and M.~Milgrom,
``Does the missing mass problem signal the breakdown of Newtonian gravity?''
\textit{Astrophys.\ J.}\ \textbf{286}, 7 (1984).

\end{thebibliography}


%=============================================================================
% APPENDICES
%=============================================================================
\newpage
\appendix

\section{Derivation of Psi-Channel Loss Formula}
\label{app:psi_loss}

Starting from the DFD action principle (Alcock 2026, Eq.~2.6), the
scalar field $\psi$ couples to the electromagnetic stress-energy tensor
through the coupling constant $\lambda$. For an SRF cavity with stored
energy $U_0$ and volume $\Vcav$, the EM energy density is
$u_{\rm EM} = U_0/\Vcav$.

The psi-field perturbation sourced by the cavity EM field:
\begin{equation}
\delta\psi_{\rm cav} = \frac{(\lambda-1)\,\Heff\,u_{\rm EM}}
{\Meff\,\Omega_\psi^2},
\end{equation}
where $\Meff = \muG\,\Omega_\psi^2/((\lambda-1)^2\,G\,\Vcav^{-1})$ is
the effective inertial mass of the psi-mode (distinct from the
gravitational mode mass $M_\psi^{\rm grav} = c^2\Vcav/(4\pi G)$; this
28-order discrepancy was resolved in W3i4).

The psi-perturbation back-reacts on the cavity, producing an additional
surface resistance:
\begin{equation}
R_\psi = \frac{(\lambda-1)^2\,\Heff^2}{Q_\psi}
\cdot\frac{4\pi G\,u_{\rm EM}}{\muG\,c^2}\cdot G_{\rm geom},
\end{equation}
where $G_{\rm geom}$ is the cavity geometrical factor. The corresponding
quality factor:
\begin{equation}
\frac{1}{Q_\psi} = \frac{R_\psi}{G_{\rm geom}} =
\frac{(\lambda-1)^2\,\Heff^2}{Q_\psi^{\rm field}}
\cdot\frac{4\pi G\,u_{\rm EM}}{\muG\,c^2}.
\end{equation}

The total cavity quality factor is:
\begin{equation}
\frac{1}{Q_0} = \frac{1}{Q_{\rm BCS}} + \frac{1}{Q_{\rm res}}
+ \frac{1}{Q_\psi},
\end{equation}
and the multi-frequency ratio diagnostic exploits the fact that
$Q_{\rm BCS}^{-1} \propto \omega^2$, $Q_{\rm res}^{-1} = \text{const}$,
and $Q_\psi^{-1} \propto \Heff^2(\omega)$, all with different frequency
dependences.


\section{Overlap Integral Computation}
\label{app:overlap}

The psi-field overlap integral $\Heff(\omega)$ for mode $n$ is defined:
\begin{equation}
\Heff^{(n)} = \frac{\int_\Vcav |\mathbf{E}_n(\mathbf{r})|^2\,
u_\psi(\mathbf{r})\,d^3r}
{\sqrt{\int_\Vcav |\mathbf{E}_n|^4\,d^3r \cdot
\int_\Vcav u_\psi^2\,d^3r}}.
\end{equation}

For the fundamental scalar mode (approximately uniform $u_\psi \approx
\text{const}$), this reduces to:
\begin{equation}
\Heff^{(n)} \approx \frac{\langle |\mathbf{E}_n|^2\rangle}
{\sqrt{\langle |\mathbf{E}_n|^4\rangle}} =
\frac{1}{\sqrt{V_{\rm eff}^{(n)}/\Vcav}},
\end{equation}
where $V_{\rm eff}^{(n)} = (\int |\mathbf{E}_n|^2)^2 /
\int |\mathbf{E}_n|^4$ is the effective mode volume.

For a TESLA 9-cell cavity (computed by finite-element analysis):
\begin{align}
\text{TM}_{010}&: \quad V_{\rm eff}/\Vcav = 0.74, \quad
  \Heff = 3.0\ee{-5}, \\
\text{TE}_{011}&: \quad V_{\rm eff}/\Vcav = 0.92, \quad
  \Heff = 1.5\ee{-5}, \\
\text{TM}_{011}&: \quad V_{\rm eff}/\Vcav = 0.85, \quad
  \Heff = 2.1\ee{-5}.
\end{align}

These values yield the ratios quoted in Table~\ref{tab:freq_scaling}.

\end{document}
